% Chapter 3

\chapter{Main Matter}% Main chapter title
\label{Chapter3}

\section{Parts}%
\label{cmd:part}

The highest form of document division is \emph{part}. This division is not
required at all. However, you can use it to give you work some structure.
There should not be more than three or four parts in the complete thesis.
Typical cases would be:
\begin{itemize}
    \item I Background or Introduction
    \item II Methodology
    \item III Application
\end{itemize}
To start a new part the command \verb!\part{title}! is used. This creates a
part page where the title is shown. It also generates an entry in the
\abb{ToC}.  The parts are automatically numbered with roman numbers but have
no influence on other numberings.

\section{Chapters}%
\label{sec:chapter}

The most important document division is the chapter. A new chapter is started
with \verb!\chapter{title}!. This creates a chapter page and generates an
entry in the \abb{ToC}. How the chapter page is designed as described in
\autoref{sub:chapterdesign}. The chapters are numbered with arabic numbers.
All following sections and subsections are in the chapter and thus, have the
chapter number.  The header also contains the chapter information. Thus a new
chapter changes the look of the header.

Maybe at different locations, the name of the chapter should be different.
There are three positions: First the entry in the \abb{ToC}. Second, the entry
in the header of the page and third is the title on the chapter page. It is
possible to give for each of the three locations a different title. Different
title version can be useful as an example when the title is too long for the
header of the page. For this the \verb!\chapter! command has two optional
arguments.  The command \verb!\chapter[TOC][HEADER]{TITLE}! generates in the
\abb{ToC} the entry \quo{TOC}. In the header, it generates the entry \quo{HEADER}
and on the chapter page the entry \quo{TITLE}.

\subsection{Mini TOC}%
\label{opt:mintoc}

At the beginning of every chapter in the main matter, a mini \abb{ToC} is
printed. This contains all sections and subsections of the chapter. The
advantage here is that the reader gets an overview directly at the beginning
of each chapter and doesn't have turn back to look at the main \abb{ToC}.

However, this needs some space and some people say it is page flay. So you may
want to disable the mini \abb{ToC}. To do this, set the class option
\clsopt{nominitoc}.\indexdef{nominitoc} When the mini \abb{ToC} is disabled
the text of the chapter begins directly on the same page under the chapter
title.


\section{Labels and References}%
\label{cmd:ref}

For internal references, the \verb!\label{key}! and the \verb!\ref{label}!
commands are normally used. This is fine and also works for this class.
However, there is a second option: the \verb!\autoref{label}! command. The
difference is that \verb!\ref! only returns the number of the referenced
object, where \verb!\autoref! also returns a suitable prefix depending on its
type (e.\,g.\ figure, section etc.). An example is shown in \autoref{alg:ref}.

The upside is that the generated hyperlinks already start at the prefix, so
they can be clicked more easily. It also reduces the possibility of
referencing the wrong object.  In rare cases, it makes it easier to change the
type of a float or a section (from/to subsection) because you do not need to
change the type prefix for every reference.

The downside is that \LaTeX{} needs to know to every type and how it should be
named.  Also, referencing more than one object is possibly not useful. Every
element that is predefined by this class can be used here. As long you do not
introduce new floats or theorems, \verb!\autoref! is the simplest way to go.
In the manual, \verb!\ref! is only used to reference more than one object (see
\autoref{sub:chapterdesign}).

A question naturally arising when using \verb!\autoref! is: \quo{How is the
type prefix displayed?} Should the first letter be uppercase or lowercase? The
default in this template is that every type prefix is spelled with an
uppercase first letter.  However, this behavior can be changed with the class
option \clsopt{smallautoref}\indexdef{smallautoref}. When it is specified, all
type prefixes are written spelled with lowercase letters.

\begin{algorithm}
    \captiont{Referencing alternatives}{%
            The two referencing methods that create the same output.
        }%
    \label{alg:ref}
    \begin{code}{latex}
See Figure~\ref{fig:veelo}.
See \autoref{fig:veelo}.
    \end{code}
\end{algorithm}

\section{Citations}%
\label{sec:cite}

There are two main citation commands that you should use: \verb!\citet! and
\verb!\citep!.  \verb!\citet! is used when the citation is read as part of the
\textbf{t}ext (\ie when using the author name as a noun), \eg when
saying that the \pkg{biblatex} package was written by \citet{web:biblatex}.
\verb!\citep! is used when the entire citation is in \textbf{p}arentheses
(which should be the default case), \eg when simply saying that we use the
\pkg{biblatex} package \citep{web:biblatex}.

Depending on whether the \clsopt{bibtex} option is used for loading the class,
citations are provided by the \pkg{natbib} package based on \pkg{bibtex}
(\clsopt{bibtex} given), or with \pkg{biblatex} in the \pkg{natbib}
compatibility mode (\pkg{bibtex} \emph{not} given). Since \pkg{biblatex}
is more modern and robust, you should not specify the \pkg{bibtex} option
unless you have a good reason to do so. By default, the \cmd{author-year}
style is used. Thus, not a number but the author name and year is printed.
This is more informative than only giving a number that you have to remember
and look up in the list of references to know the author.  Though numbered
citations are more common in papers, you should consider that your thesis is
going to be a much bigger document and the bibliography may become huge.

\section{Margin Notes}%
\label{cmd:margin}

There is the possibility to print notes in the margin of the space. More
precisely on the outside margin of the page. These notes are very prominent
and seen easily but destroy the reading flow because they are outside of the
normal reading area. This can be used for information that should be easily
found but should not be used too often because then it loses its importance
and only disturbs the reading flow.

To do such a margin note use the command \verb!\marginnote{text}!. An example
is in this manual the class option. Every time the class options are presented
in different sections all over the manual, the command appears as a margin
note.  Thus, every time a new class option is introduced a margin note gives a
hint that it is presented here. This makes it possible to scan quickly for a
specific class option.

\section{Definition Index}%
\label{cmd:defindex}

One of the most popular uses of margin notes is to give a reader the position
of a definition. Then when a term is used the reader can easily find it and
check what it means. This concept can be extended by creating an index at the
end of the thesis where for every definition a  site is given where to find
it. Such that from anywhere in a large thesis fast the right position of a
definition can be found.

This can be archived with a special command. The template gives the
\verb!\indexdef{name}! command. This command creates a margin note with the
given name at the side. Further when the
\clsopt{definitionindex}\indexdef{definitionindex} class option is given an
entry in the Definition Index is created.

However, not every time you want to create an index you also want to create a
margin note. Thus, there are two different commands also given in to create
index entries. First the \verb!\indexnomn{name}! command. This command creates
an entry in the index but does not print anything at the position. It could be
combined with the \verb!\marginnote{text}! command to create different content
in the margin note and the index.

The second other command is \verb!\indexbf{name}!. This command also creates
an index entry and print \quo{name} in bold letters.  Thus, a highlighting of the
position is possible without the disturbing character of margin notes. It is
particularly useful when small margins are used.

\section{Abbreviations}%
\label{sec:abb}

% Abbreviation definitions look like this:
%   \defabb{UTC}{Coordinated Universal Time}
% But you should add them to misc/custom_defs.tex, not to the chapter content!
Every thesis uses abbreviations. An abbreviation can be a generally used term
like time zones (\abb{UTC}) or a topic-specific thing (\abb{DAG}) but also
very specific problem descriptions (\abb{mFAS} Problem). So it is necessary to
keep track of these abbreviations in a List of Abbreviations.  To add a new
abbreviation use the \verb|\defabb{short}{long}| command.

Besides this list, other things must be considered. If an abbreviation is
first introduced a long version must be given to explain the abbreviation. In
the following, only the short version is used. However, after some pages of
not using the abbreviation, it is useful to use the long version again so that
the reader gets a reminder what this abbreviation stands for.  Another thing
that must be considered is that sometimes the plural (\abbs{DAG}) of an
abbreviation is used (per default an \quo{s} is added to the end).

So the abbreviation engine gives more than the \quo{List of Abbreviations}.
It also gives the possibility to access earlier defined abbreviations with the
\verb!\abb! (for plural \verb!\abbs!) command. This returns the long version
with the short version in parentheses after the long version was not used on
the pages before. An example of the syntax can be found in \autoref{alg:abb}.

The default plural form with simply appending an \quo{s} is not always the
correct way to go (\abbs{SV}). For such situations, the alternate command form
\verb|\defabb[label]{short}[plural short]{long}[plural long]|
with the optional arguments \verb|plural short| and \verb|plural long| can be
used.  One is the short version in plural and the other the long version in
the plural. Note that it is possible to give only one of the two plural forms,
the other is then simply created by adding \quo{s}. The \verb|label| argument
allows to assign a label to the abbreviation such that it can be referred to
using \verb|\abb{label}|, which then still prints the supplied
short or long form.  With the two-argument command, the short form of the
abbreviation is itself used as the label. However, for complex abbreviations
including formatting and other commands, this is not possible.

\begin{algorithm}
    \captiont{The abbreviations engine.}{%
        An example definition of Abbreviation and the use of it. It is also an
        entry to the List of Abbreviations created.
    }%
    \label{alg:abb}
    \begin{code}{latex}
\defabb{DAG}{Directed Acyclic Graph}
The life is a \abb{DAG}.
%The life is a Directed Acyclic Graph (DAG).
Many \abbs{DAG} together create a \abb{DAG}.
%Many DAGs together create a DAG.
\defabb{SV}[SVs]{Super Vertex}[Super Vertices]
\abbs{SV}
%Super Vertices (SVs)
    \end{code}
\end{algorithm}

The default value of pages between two uses of \verb!\abb! before the long
version is printed again is 3. This can be changed with the \clsopt{maxdis}
class option.\indexdef{maxdis} A value of 0 means that only the first
appearance is printed with the long version all other occurrences have the
short version.

The functionality described here comes from the glossary package. It is a
large package. Thus, it makes sense not load it when you do not want any
\quo{List of Abbreviation} at all. This deactivation is done with the class
option \clsopt{nostatistics}.\indexdef{nostatistics} After the deactivation,
the abbreviation commands still exists but do different things. The
\verb!\defabb! command does nothing, the \verb!\abb! command simple print the
abbreviation, and the \verb!\abbs! command print the abbreviation with an
\quo{s} at the end. This behavior makes it possible to use them and then later
activate the glossary package again and have the full functionality with the
commands.

\section{Symbols}%
\label{sec:sym}

In the back matter, a List of Symbols is presented. To make this possible,
symbols need to be defined first.  To introduce a new symbol, use the
command \verb!\sym{symbol}{Description}!.  The symbol can then be used in a
mathematical environments as well as in text. It is also added to the List of
Symbols using the provided description.  An example is shown in
\autoref{alg:sym}. While symbols may be defined within the document body, \eg
directly where they are used for the first time, it may be more manageable to
maintain all definitions in \file{misc/custom\_defs.tex} to have them all in a
single place.

\begin{algorithm}
    \captiont{The sym command.}{%
        An example definition of Symbol. The result can be seen in the List of
        Symbols.
    }%
    \label{alg:sym}
\begin{code}{latex}
\sym{C}{The circumference of a cycle.}
\sym{\pi}[kreis]{Ratio of a circle's
    circumference to its diameter.}
We use the new cool symbol in math
$\symkreis = 3.14\dots$ and in text \symkreis.

\operatorsym{lca}[lca]{The last common ancestor.}
Let \symlca be the last common ancestor. When $v$ has
two children $u$ and $w$ then is $\symlca(u,w)=v$.
\end{code}
\end{algorithm}

This is some extra work for the \quo{List of Symbols}. However, this list
really helps the reader to understand complex formal definitions. It is also
very helpful for the writer. The list makes it easy to determine if a symbol
is used twice with varying meaning. This should be avoided, even in such a
complex work as a thesis.

The syntax is simple. The definition has an optional field in the middle that
defines a name: \verb!\sym{symbol}[name]{Description}!. This name can then be
used to reference back to the symbol using the command \verb!\sym*name*!. If
the name is omitted, the symbol itself is used as the name. This works only if
the symbol consists of regular letters, otherwise a name must be provided.

Sometimes you may want to define a custom math operator as a symbol. This can
be done by using the \verb!\DeclareMathOperator{\lca}{lca}! command in
\file{misc/custom\_defs.tex} (or another place in the preamble).  Then, you
can create a symbol using this operator. However, if you really want, you can
also define an operator symbol in the document body with the help of
the \verb!\operatorsym{symbol}[name]{Description}! command. Note that this
makes it necessary to use the \verb!\sym*name*! command to refer to the
operator. If you do not wish to have the \quo{sym} prefix, you need to declare
the math operator using \verb!\DeclareMathOperator{\lca}{lca}!.

The symbol engine has a problem with the \quo{$\vert$} symbol. However, this
can be fixed by not using direct \quo{$|$} but one of the two command
replacements: \verb|\vert| or \verb|\mid|. The first one is more or less the
same as \quo{$|$} and is used for example to get the absolute value of a number.
The second have more space around it and is used in constructive set
definitions.

\operatorsym{lca}[lca]{The last common ancestor.}[T]
\sym{\pi}[kreis]{The ratio of a circle's circumference to its diameter.}[C]
\sym{C}{The circumference of a cycle.}[C]

\subsection{Ordering and Categories}
The ordering of the symbols fellow normal the ASCII code
(+<1<:<A<\textbackslash<a). Note symbols define with \verb!\operatorsym! are
placed with \verb!\operator! at the beginning. Especially the fact that
\quo{\textbackslash} is placed between \quo{A-Z} and \quo{a-z} in the ordering is
often some way counterintuitive. To improve on this, you can customize the
order.

Both \verb!\sym! and \verb!\operatorsym! have an optional argument at the
end. This defines the ordering. It is important to note that every symbol for
which this last argument is given is placed behind every symbol without a last
argument.  Thus, it often only makes sense to give such an ordering argument
either to all symbols or to none. As an example if \quo{\symlca} should be
placed as lca in the ordering the definition should look like:
\verb!\operatorsym{lca}{The last common ancestor.}[lca]!.

Some times you want to group you symbols. This can be archived in two steps.
First you must define the groups in the preamble. This is done with
\verb!\addSymbolCategory{C}{Cycle}! command. The first argument must be a
\emph{capital letter} and the second argument the name of the category.
Second, first argument can then be used as last argument on the symbol
definition to get a element in the category. It is important to note that this
really only work with one capital letter. When more then one letter is given or
a small letter \LaTeX{} can not use the category. An example can be seen in
\autoref{alg:symcat}.

\begin{algorithm}
    \captiont{The symbol categories.}{%
        An example definition of Symbol categories. The result can be seen in
        the List of Symbols.
    }%
    \label{alg:symcat}
\begin{code}{latex}
%Categories can only defined in the preamble.
\addSymbolCategory{C}{Cycle}
\addSymbolCategory{T}{Tree}

\begin{document}
% Symbols can defined in the document body.
\operatorsym{lca}[lca]{The last common ancestor.}[T]
\sym{\pi}[kreis]{The ratio of a circle's
    circumference to its diameter.}[C]
\sym{C}{The circumference of a cycle.}[C]
\end{document}

\end{code}
\end{algorithm}

\subsection{Symbols with Parameters}
In many situations, you do not only have a symbol but also the symbols depend
on other things. An example would be the vertex set of a graph. There are at
least two different used notations for this $G.V$ or $V(G)$.  However, it
makes only sense when the graph (in this case $G$) is given. Thus a vertex set
symbol needs a place that can be changed.

This behavior is reached with the \verb!\macrosym! command. This command
defines a macro that can be used in the thesis and also generates an entry in
the list of symbols. Because this needs a name for the macro, its signature is
different than the other methods.

The command \verb!\macrosym{content}{name}[C]{d}[p1][p2]...!
creates the macro that is \verb!\sym*name*!. The content is what is printed.
It is given in latex macro form which means for the first parameter is \quo{\#1}
and the second \quo{\#2} and so on. The category (C) and description (D) are the
same as by the other commands. The command ends with a list of the parameters
that are used in the index. There are up to 5 parameters possible. When you
want a macro with four parameters you need to give four at the end of the
command.

\begin{algorithm}
    \captiont{The macro symbol command.}{%
        Some examples of macro symbols and how they can be used.
    }%
    \label{alg:symmac}
\begin{code}{latex}
%Let G be the 'graph' category.
\macrosym{V(#1)}{vertexset}[G]{The vertex set of $G$}[G]
\macrosym{E(#1)}{edgeset}[G]{The vertex edge of $G$}[G]
\macrosym{#1[#2]}{indsub}[G]{Creates induced subgraph of $G$
 with vertex set $U$ }[S][U]
%Let S be the 'set' category.
\macrosym{###1}{cardinality}[S]{The number of elements
in set $U$}[U]

When $D$ is a directed graph and $U\subset\symvertexset{D}$ and
$H=\symindsub{D}{U}$ have no cycles then is
$\symcardinality{\symedgeset{H}}\leq\frac{1}
{\symcardinality{U}(\symcardinality{U}-1)}$.
\end{code}
\end{algorithm}
%Let G be the 'graph' category.
\macrosym{V(#1)}{vertexset}[G]{The vertex set of $G$}[G]
\macrosym{E(#1)}{edgeset}[G]{The vertex edge of $G$}[G]
\macrosym{#1[#2]}{indsub}[G]{%
    Creates induced subgraph of $G$ with vertex set $U$ }[S][U]
%Let S be the 'set' category.
\macrosym{\##1}{cardinality}[S]{The number of elements in set $U$}[U]

\autoref{alg:symmac} creates this output:
\begin{quote}
    When $D$ is a directed graph and $U\subset\symvertexset{D}$ and
    $H=\symindsub{D}{U}$ have no cycles then is
    $\symcardinality{\symedgeset{H}}\leq
    \frac{1}{\symcardinality{U}(\symcardinality{U}-1)}$.
\end{quote}

\section{Floats}%
\label{sec:float}

There are three basic types of floats: figure, table and algorithm. Beside
this, there are two special types: SCfigure and SCtable. These are the same as
the floats without \quo{SC}. The difference is that in these floats the caption
is placed beside the figure (table). This can be used with narrow figures
(tables) to save space. An example of SCtable can be seen at
\autoref{tab:fontoption}.

\subsection{Figures}

The include \verb!\includegraphics! command of a figure needs the path to the
figure. The important part is that this path must be given relative to the
main.tex file. For an example is the picture veelo.pdf in the Figures folder
linked as \quo{graphics/ell.pdf} in every Chapter.


\subsection{Captions}%
\label{sub:cap}

The caption can be used with the normal \verb!\caption! command. Remember,
that the optional argument that sets the title on the \quo{List of \dots} page.
However, an alternative in this class is the \verb!\captiont! command that has
a second mandatory argument that sets the title and also writes it as bold at
the beginning of the caption. An example is shown in \autoref{alg:caption}.

All captions use a little smaller font than the rest of the text. This helps
to distinguish between text and caption of a float. This supports the reading
flow.

The caption is placed above or under the content of the float depending on the
position of the caption command. However, normally the figure caption should
be placed under the content, in a table above the content and in an algorithm
also above the content.

\begin{algorithm}
    \caption{%
        An example of how the \cmd{captiont} command can be used.
    }%
    \label{alg:caption}
    \begin{code}{latex}
\begin{SCtable}
    \captiont{Font values.}%
    {The values that the class option \quo{font} can have.}%
    \label{tab:fontoption}
    \begin{tabular}{l r}
        \toprule
        \tblhead{Value} & \tblhead{Used Font}\\
        \midrule
        lm  & Latin modern without serifs (default)\\
        lms & Latin modern with serifs\\
        cm  & Computer modern without serifs\\
        cms & Computer modern with serifs\\
        \bottomrule
    \end{tabular}
\end{SCtable}
    \end{code}
\end{algorithm}

\section{Code}%
\label{sec:code}

This template provides two different ways to typeset program code. One is
meant for code written in a real programming language and includes syntax
highlighting, the other is for pseudo-code that is more abstract than real
code.  These variants use the \cmd{code} and the \cmd{algorithmic}
environments, respectively.  However, in both cases, a surrounding
\emph{float}, which is automatically positioned and provides a caption that
can be labelled and referenced, is useful. This feature is available in form
of the \cmd{algorithm} environment, which then includes either a \cmd{code} or
a \cmd{algorithmic} environment. An example is shown in \autoref{alg:algo}.

\begin{algorithm}
    \caption{%
        Example of an \cmd{algorithm} float enclosing a
        \cmd{code} environment to display its own code.
    }%
    \label{alg:algo}
    \begin{code}{latex}
\begin{algorithm}
    \caption{%
        Example of an \cmd{algorithm} float enclosing a
        \cmd{code} environment to display its own code.
     }%
    \label{alg:algo}
    *code environment*
\end{algorithm}
    \end{code}
\end{algorithm}

\subsection{Pseudo Code}
To typeset pseudo code, the \pkg{algorithmic} package is used. See
\autoref{algo:algorithmic} for how the code looks like
For more details what commands exist, please have a look at the
manual.\footnote{\url{ftp://ftp.fu-berlin.de/tex/CTAN/macros/latex/contrib/algorithms/algorithms.pdf}}

\begin{algorithm}
    \caption{%
        The \cmd{algorithmic} environment can be used to typeset
        pseudo-code.
    }%
    \label{algo:algorithmic}
    \begin{code}{latex}
\begin{algorithmic}
    *pseudo-code*
\end{algorithmic}
    \end{code}
\end{algorithm}

\subsection{Real Code with Highlighting}%
\label{sub:codehigh}

The \cmd{code} environment is used to typeset code written in a real
programming language. Passing the name of a supported programming language,
\eg \verb!\begin{code}{python}!, allows the package to apply syntax
highlighting to the code. An example is given in \autoref{alg:code}.  Though
\cmd{code} environments can be used by itself to place the code exactly at the
given position, you should usually enclose it by an \cmd{algorithm}
environment to allow automatic positioning and provide a proper caption as in
\autoref{alg:algo}.

For code highlighting, two different packages are used. The default package is
the \pkg{listings} package. This is nice because it works out of the box on
any \LaTeX{} distribution. To achieve nicer syntax highlighting (with color),
another package is used: the \pkg{minted} package. It depends on the external
\prog{Python} program \prog{Pygments}. In order to use it, you need to install
them. More information on this can be found
online.\footnote{\url{https://github.com/gpoore/minted}} The manual is created
with \pkg{minted}.  To use \pkg{minted} instead of \pkg{listings} you set the
class option \pkg{minted}.\indexdef{minted}

However, which package you use does not have an influence on the environment
name, which is always \cmd{code}. This makes it possible to switch between
both options on the fly.

\begin{algorithm}
    \caption{An example how code highlighting works.}%
    \label{alg:code}
    \begin{code2}{latex}
\begin{code}{python}
    *code*
\end{code}
    \end{code2}
\end{algorithm}


\section{Math Environments}%
\label{sec:theo}

There are six different environments that belong to this class. They all have
their own uses. The explanations of them are taken from \citet{web:theorems}.


\subsection{Proposition}
A proved and often interesting result, but generally less important than a
theorem.
\begin{code}{latex}
\begin{proposition}
    *text*
\end{proposition}
\end{code}

\subsection{Definition}
A precise and unambiguous description of the meaning of a mathematical term.
It characterizes the meaning of a word by giving all the properties and only
those properties that must be true.
\begin{code}{latex}
\begin{definition}
    *text*
\end{definition}
\end{code}

\subsection{Corollary}
A result in which the (usually short) proof relies heavily on a given theorem
(we often say that “this is a corollary of Theorem A”).
\begin{code}{latex}
\begin{corollary}
    *text*
\end{corollary}
\end{code}

\subsection{Theorem}
A mathematical statement that is proved using rigorous mathematical reasoning.
In a mathematical paper, the term theorem is often reserved for the most
important results.
\begin{code}{latex}
\begin{theorem}
    *text*
\end{theorem}
\end{code}

\subsection{Lemma}
A minor result whose sole purpose is to help in proving a theorem.  It is a
stepping stone on the path to proving a theorem.
\begin{code}{latex}
\begin{lemma}
    *text*
\end{lemma}
\end{code}

\subsection{Proof}
Every Corollary, Theorem and Lemma must be proven. Thus, such an environment,
a Proof environment should follow.
\begin{code}{latex}
\begin{proof}
    *text*
\end{proof}
\end{code}

\section{Todo Notes}%
\label{sec:todo}

Often, some task is left undone and has to be finished later.  This should be
made visible in the thesis. The \verb!\todo{}! command is the solution to this
problem.  It creates a colored note containing the todo text.\todo{This is an
example of a todo note. Try changing the \clsopt{todo} option of the class and
see what happens.} The exact behaviour can be controlled with the class option
\clsopt{todo}, which can be set to any of the three values \emph{margin},
\emph{plain}, and \emph{off}.  The \emph{margin} setting creates orange boxes
located in the margin, keeping the layout intact and separating text and
tasks. Additionally, a list of all remaining tasks is generated and placed on
the first page of the thesis. This is probably the most elegant solution, but
relies on the external package \pkg{marginnote}.  The \emph{plain} option
simply inserts the task as red text, enclosed in brackets and with a smaller
font size than the main text. It has no external dependencies.  The value
\emph{off}, finally, allows to temporarily disable the output of todo notes
completely without deleting them.  This is useful for creating drafts of the
thesis that should not contain the note text.

% End of chapters/chapter_3.tex
